\documentclass[10pt,a4paper]{article}
\usepackage[margin=22mm]{geometry}
\usepackage{amsmath,booktabs,hyperref}
\hypersetup{colorlinks=true,urlcolor=blue,linkcolor=blue,citecolor=blue}
% Generated by manuscript.py; do not edit.
\newcommand{\ConservativeFront}{0.213862}
\newcommand{\ConservativeEnergy}{6.433e-13}
\newcommand{\FrozenFront}{1.542954}
\newcommand{\FrozenEnergy}{4.248e-02}
\newcommand{\MaxResidual}{2.977e-13}
\newcommand{\Similarity}{0.3064239054}
\newcommand{\HeatAlgebraic}{2.842e-13}
\newcommand{\BaselineTemperature}{0.000e+00}
\newcommand{\BaselineEnergy}{1.393e-14}
\newcommand{\ConservativeMeshOneEarly}{0.699333}
\newcommand{\FrozenMeshOneEarly}{0.177702}
\newcommand{\CoarseTimeEarly}{0.030316}
\newcommand{\FineTimeEarly}{0.066729}

\title{Energy Conservation and Front Accuracy\\Across a Regularized Melting Interval}
\author{}
\date{Transport Verification Lab: November 2026}
\begin{document}
\maketitle
\begin{abstract}
When a time step crosses phase endpoints, does an accurate linear solve also
conserve the material's true energy? We connect a conservative enthalpy solver,
an intentionally frozen apparent-capacity comparator, independent heat and
Stefan references, and Python/C++ implementations in a reproducible verification
study. On eleven fixed front settings, the conservative cumulative scaled
energy defect stays below $7\times10^{-13}$. Both methods meet the prescribed
2\% front-error target at the finest setting. Unfavorable comparisons and
nonmonotone errors are retained: conservation does not imply uniformly superior
front accuracy. Independent smooth heat tests separate first-order temporal and
second-order spatial error. The contribution is an original connected
verification study of established methods, not a new discretization or
experimental physical validation.
\end{abstract}
\section{Question, model and numerical connection}
Enthalpy formulations for phase change are established~\cite{voller}. Here we
ask how replacing a finite enthalpy increment by an old-state tangent affects
energy accounting and front accuracy when a step crosses a melting interval.
DTU's Modelling in Materials and Manufacturing Engineering (MMME) topics on
solidification and implicit conduction motivate this study. Equations, code
and synthetic experiments are original; no private course files, exercise
solutions or teaching code are distributed. The six imported Transport modules
remain versioned snapshots. The new \texttt{phasechange} package and C++17
kernels form a connected study within that laboratory.

Let $\rho,c,k,L$ denote constant density, specific heat, conductivity and
latent heat per unit mass. On a stationary one-dimensional slab of area $A$,
define volumetric enthalpy $H$ and liquid fraction $f$:
\begin{equation}
H(T)=\rho c(T-T_{\rm ref})+\rho Lf(T),\qquad
f(T)=\min\{1,\max[0,(T-T_s)/w]\},\quad w=T_l-T_s>0.
\end{equation}
The model is $\partial_tH=k\partial_{xx}T+q$, with volumetric power $q$.
There is no advection, density change, segregation or deformation.
The constitutive inverse is exact on each branch. Clipping in $f$ defines the
material law; accepted temperatures are never clipped to hide errors and must
be nonnegative kelvin.

For cell volume $V=A\Delta x$, internal conductance $g=kA/\Delta x$ and
prescribed-wall conductance $2g$, conservative backward Euler solves
\begin{equation}\label{eq:step}
R_i=V[H(T_i^{n+1})-H(T_i^n)]
 +\Delta t\sum_jg_{ij}(T_i^{n+1}-T_j^{n+1})-\Delta tVq_i=0.
\end{equation}
The sum includes a prescribed wall as a fixed neighbor; insulation omits that
face. Newton steps use a right-sided tangent and a positive definite
tridiagonal Jacobian. A scalar line search zeros the directional derivative
of the strictly convex residual potential when a full step passes its minimum.
The stopping test is $\|R\|_\infty/[V(\rho cw+\rho L)]\le10^{-10}$;
iteration exhaustion is an error. The C++17 backend supplies residual,
Jacobian and $LDL^T$ solves through an explicit Python interface. Both backends
share the Python safeguard and ledger; agreement checks implementation, not
independent validity of that algorithm.

The frozen comparator replaces storage in (\ref{eq:step}) by
$VH'_+(T_i^n)(T_i^{n+1}-T_i^n)$ for the entire step, including cooling from a
phase endpoint. This deliberately uniterated approximation is not a claim
about all apparent-capacity methods. Its linear residual can be small while
the true enthalpy balance fails. With $Q_j$ the signed boundary/source heat
entering the slab in step $j$, define
\begin{equation}
D_n=\frac{|\sum_iV[H(T_i^n)-H(T_i^0)]-\sum_{j=1}^nQ_j|}
{A\ell(\rho cw+\rho L)},\qquad D_{\max}=\max_nD_n.
\end{equation}
Reporting the largest transient defect avoids hiding drift through final-cycle
cancellation. Algebraic closure and temperature/front accuracy are distinct.

\newpage
\section{Independent references and experimental protocol}
Uniform heating and cooling remove spatial error. An independent
event-traversal reference computes times to each phase endpoint and integrates
with the appropriate constant capacity; it calls neither the enthalpy inverse
nor the spatial solver. Six schedules with 1, 2, 4, 8 and 16 subdivisions
give baseline temperature error scaled by the 2 K interval and energy
error scaled by $\rho cw+\rho L$ of at most
\texttt{\BaselineTemperature} and \texttt{\BaselineEnergy}.
Constitutive inverse, monotonicity and endpoint tests complement this check.

For $L=0$, an insulated slab has exact cell-average cosine amplitudes
\begin{equation}
a_i=\frac{\sin(\pi/(2N))}{\pi/(2N)}\cos(\pi x_i/\ell),\qquad
\mu=\alpha(\pi/\ell)^2,\qquad
\mu_h=\frac{4\alpha}{\Delta x^2}\sin^2\frac{\pi}{2N},
\end{equation}
where $\alpha=k/(\rho c)$. Continuum and semidiscrete temperatures are
$310+8a_i e^{-\mu t}$ and $310+8a_i e^{-\mu_ht}$ K. Backward Euler
multiplies the discrete amplitude by $(1+\mu_h\Delta t)^{-1}$ per step.
Comparing numerical steps with exact semidiscrete evolution isolates temporal
error; comparing the two exact exponentials isolates spatial error without
time contamination. An independently assembled incidence matrix and dense
matrix exponential check the discrete oracle; quadrature checks cell averages.

Synthetic constants are $\rho=1200$ kg/m$^3$, $c=900$ J/(kg K),
$k=1.8$ W/(m K), and $L=90000$ J/kg when phase change is active.
Heat studies use $\ell=0.1$ m and $t=1/\mu=607.927101854$ s: 32 cells with
20--160 steps for time error, and 10--80 cells with exact discrete evolution
for spatial error. These heat settings were specified on day 5; the Stefan
protocol below was fixed on day 1. Table entries are maximum cell-temperature
errors; successive order is $\log_2(E_{\rm coarse}/E_{\rm fine})$.
% Generated by manuscript.py; do not edit.
\begin{center}\small
\begin{tabular}{rrrrrr}\toprule
Steps & Time error (K) & Order & Cells & Space error (K) & Order\\\midrule
20 & 7.1907e-02 & -- & 10 & 2.3829e-02 & --\\
40 & 3.6322e-02 & 0.9853 & 20 & 6.0278e-03 & 1.9830\\
80 & 1.8255e-02 & 0.9926 & 40 & 1.5114e-03 & 1.9958\\
160 & 9.1510e-03 & 0.9963 & 80 & 3.7812e-04 & 1.9989\\
\bottomrule\end{tabular}\end{center}

The largest error relative to the algebraic backward-Euler amplitude is
\texttt{\HeatAlgebraic} K. These orders concern smooth constant-property
heat conduction, not a moving phase front.

For a hot wall $T_w$ and initially solid material at $T_m$, the independent
one-phase sharp Stefan reference is
\begin{align}
s(t)&=2\lambda\sqrt{\alpha t},&
\lambda e^{\lambda^2}\operatorname{erf}(\lambda)
 &=\frac{c(T_w-T_m)}{L\sqrt{\pi}},\\
T(x,t)&=T_w-(T_w-T_m)
\frac{\operatorname{erf}(x/(2\sqrt{\alpha t}))}{\operatorname{erf}(\lambda)}
\quad(x<s),&
T(x,t)&=T_m\quad(x\ge s).
\end{align}
The constant-conductivity, prescribed-wall limit of the similarity
construction is discussed in~\cite{ceretani}; that work treats solidification
and a broader boundary model. The melting sign convention here is checked
independently through the heat PDE, boundary temperatures and Stefan flux jump.
Bisection of a monotone logarithmic equation gives $\lambda=\Similarity$.

The fixed front protocol uses $T_m=310$ K, $T_w=330$ K, $A=1$ m$^2$ and an
insulated far wall at $\ell=0.1$ m. Regularization has $T_s=T_m$, $T_l=T_m+w$
and zero initial liquid fraction. Observations at 30, 60, 90 and 120 s exclude
the singular initial instant. The numerical front is the interpolated
decreasing $f=0.5$ crossing between cell centers; unresolved crossings are
explicit failures. Four-level studies vary $N=100,200,400,800$,
$\Delta t=0.5,0.25,0.125,0.0625$ s, or $w=1,0.5,0.25,0.125$ K, holding the
other two at their finest values. Doubling $\ell$ and $N$ preserves $\Delta x$.
Shared finest cases are computed once: eleven settings, each with both
methods and both backends.

\newpage
\section{Measured endpoint, energy and front behavior}
Superscripts $H,F$ label conservative enthalpy and frozen capacity.
Uniform $E_T$ is the largest error against the event reference. The interval
is 309--311 K and the volume is 1 m$^3$. Solid heating starts at 305 K for
1 s; mushy heating/cooling start at 309.5/310.5 K for 2 s. Power is
$\pm2.7$ MW/m$^3$. Crossing tests start at 305 or 330 K with the same signed
power for 50 s. The closed cycle heats from 305 K for 50 s and cools for
another 50 s. This table
uses one step per segment; all subdivision histories remain in the JSON.
Scientific notation uses \texttt{e} for powers of ten.
% Generated by manuscript.py; do not edit.
\begin{center}\small
\begin{tabular}{lrrrr}\toprule
One step per segment & $E_T^H$ (K) & $E_T^F$ (K) & $D^F_{\max}$ & $D^F_{\rm end}$\\\midrule
solid heating & 0.00e+00 & 0.00e+00 & 0.000e+00 & 0.000e+00\\
mushy heating & 0.00e+00 & 0.00e+00 & 1.226e-14 & 1.226e-14\\
mushy cooling & 0.00e+00 & 0.00e+00 & 1.226e-14 & 1.226e-14\\
cross both heating & 0.00e+00 & 1.00e+02 & 9.804e-01 & 9.804e-01\\
cross both cooling & 0.00e+00 & 1.00e+02 & 9.804e-01 & 9.804e-01\\
heat cool cycle & 0.00e+00 & 1.00e+02 & 9.804e-01 & 0.000e+00\\
\bottomrule\end{tabular}\end{center}

The 100 K error is missed latent heat divided by sensible capacity, $L/c$.
Final-cycle cancellation hides transient drift. Within a constitutive branch
the methods agree. These crossing failures concern the specified uniterated
approximation and time steps, not all capacity formulations.

A connected slab test uses $\ell=0.012$ m, $A=0.003$ m$^2$, 24 cells,
305 K initial temperature and an insulated right wall. Five 10 s base steps
use a 335 K left wall and $q=2.7$ MW/m$^3$; five more use 295 K and negative
power of the same magnitude. Each base step is subdivided uniformly.
The temperature difference is between methods, not continuum error.
% Generated by manuscript.py; do not edit.
\begin{center}\small
\begin{tabular}{rrrr}\toprule
Subdivisions & $\max|T_H-T_F|$ (K) & $D^H_{\max}$ & $D^F_{\max}$\\\midrule
1 & 69.453 & 2.451e-15 & 9.804e-01\\
2 & 73.183 & 4.171e-15 & 9.804e-01\\
4 & 74.972 & 1.399e-14 & 9.804e-01\\
8 & 68.604 & 9.861e-15 & 8.488e-01\\
16 & 34.821 & 1.961e-14 & 4.421e-01\\
\bottomrule\end{tabular}\end{center}

The nonmonotone differences show why implementation agreement cannot determine
accuracy. Ledgers retain all steps, including unsampled front-study steps.
For that separate study, define
\[
E_s=100\max_{t\in\{30,60,90,120\}}|s_h(t)/s(t)-1|.
\]
The shared finest case is labeled refined; numbered cases use the levels
listed above. The doubled-domain case uses $\ell=0.2$ m.
% Generated by manuscript.py; do not edit.
\begin{center}\small
\begin{tabular}{lrrrrr}\toprule
Setting & Cells & $\Delta t$ (s) & $w$ (K) & $E_s^H$ (\%) & $E_s^F$ (\%)\\\midrule
refined & 800 & 0.0625 & 0.125 & 0.213862 & 1.542954\\
mesh-0 & 100 & 0.0625 & 0.125 & 2.861279 & 3.068805\\
mesh-1 & 200 & 0.0625 & 0.125 & 0.839380 & 1.032510\\
mesh-2 & 400 & 0.0625 & 0.125 & 0.625741 & 1.311259\\
time-0 & 800 & 0.5 & 0.125 & 0.217875 & 11.078534\\
time-1 & 800 & 0.25 & 0.125 & 0.214978 & 7.142458\\
time-2 & 800 & 0.125 & 0.125 & 0.214434 & 2.409037\\
width-0 & 800 & 0.0625 & 1 & 1.583802 & 2.223497\\
width-1 & 800 & 0.0625 & 0.5 & 0.898116 & 1.511194\\
width-2 & 800 & 0.0625 & 0.25 & 0.277012 & 1.551998\\
double-domain & 1600 & 0.0625 & 0.125 & 0.213862 & 1.542954\\
\bottomrule\end{tabular}\end{center}

Both refined methods meet the 2\% target: $E_s^H=\ConservativeFront\%$ and
$E_s^F=\FrozenFront\%$. Across these settings the largest conservative and
frozen $D_{\max}$ are \texttt{\ConservativeEnergy} and \texttt{\FrozenEnergy}.
The largest equation residual for either method is \texttt{\MaxResidual}
on its cell storage scale. A small linear residual does not imply true-energy
closure; conservative closure does not remove discretization or regularization.

\newpage
\section{Interpretation, reproducibility and limits}
\paragraph{Retained counterexamples.}
The finest conservative maximum front error is smaller, but this is not a
universal ranking. At 30 s with 200 cells, frozen error is
\FrozenMeshOneEarly\% versus \ConservativeMeshOneEarly\% for enthalpy.
On 800 cells with $w=0.125$ K, reducing $\Delta t$ from 0.5 to 0.0625 s
increases early conservative error from \CoarseTimeEarly\% to \FineTimeEarly\%.
This is consistent with error cancellation, not evidence that coarse steps
are generally preferable. No Stefan convergence order is asserted.
Shrinking $w$ at fixed resolution can eventually underresolve the phase
interval; the width study is a sensitivity experiment, not a complete
asymptotic limit.

\paragraph{Independent evidence and its limits.}
Event tests verify local accounting. Python/C++ residual/state comparisons,
finite-difference derivatives away from kinks, one-sided endpoint checks
and native invariants assess implementation consistency. Exact discrete
heat evolution separates temporal from smooth spatial continuum error.
Sharp Stefan discrepancies include finite-interval regularization and front
interpolation as well as discretization. These factors are varied separately;
an additive error decomposition is not claimed. Doubling the domain produces
zero measured front and common-cell temperature differences for both methods
at the four observations. This supports insensitivity at floating-point
resolution in this window only, not a bound for later times or other inputs.
Residuals are below the stopping threshold; no tolerance sweep was performed.

\paragraph{Reproduction and evidence chain.}
\begin{sloppypar}
Source and synthetic results are public~\cite{code}. Use Python 3.11 or later
(CI: 3.12), NumPy/SciPy and a C++17 compiler. In a fresh virtual environment,
run \texttt{python suite.py} with \texttt{install}, \texttt{verify},
\texttt{native} and \texttt{test} in turn. Reproduce the baseline with
\texttt{python -m phasechange}, the comparison with
\texttt{python -m phasechange --study comparison}, and refinement with
\texttt{python -m phasechange.refinement}.
The comparison and refinement commands accept \texttt{--backend native} and
\texttt{--check-against} with a saved JSON path. Use distinct \texttt{--output}
paths to preserve reference data. Linux and Windows CI repeat both complete
campaigns, storing sampled fields, signed heat histories, front status,
acceptance flags and numerical source/protocol SHA-256 hashes. Tests fail on
missing evidence or numerical disagreement. The six imported snapshots remain
separately installable and hash-checked; their prior synthesis remains at the
v0.1.0 tag. Original histories are preserved.
\end{sloppypar}

After reproducing changed experiments, run \texttt{python manuscript.py generate}.
It refuses stale numerical source hashes and generates all five table/number
fragments. \texttt{python manuscript.py check} rejects stale tables.
Build with \texttt{python suite.py paper --engine tectonic} or pdfLaTeX;
the build rejects unresolved references, overfull boxes and more than five
pages including references. \texttt{python manuscript.py check-report}
checks page count and the complete source/data/table/PDF hash set.
Mutation tests cover source drift, altered measurements, missing observations,
failed gates, changed PDFs and incomplete manifests. CI verifies the committed
artifact before rebuilding it independently. Local and CI pages are rendered
for visual review, with actual outcomes in dated validation records.

\paragraph{Conclusion and scope.}
Finite enthalpy increments preserve energy across phase endpoints here;
conservation cannot replace independent accuracy checks. The model excludes
convection, segregation, unequal phase properties and hysteresis.
No calibration, experimental validation, speedup or sanitizer coverage is
claimed. Release remains conditional on all acceptance gates.

\begin{thebibliography}{9}\small
\bibitem{voller} V. R. Voller, M. Cross and N. C. Markatos.
An enthalpy method for convection/diffusion phase change.
\emph{International Journal for Numerical Methods in Engineering}
24(1), 271--284 (1987). \url{https://doi.org/10.1002/nme.1620240119}.
\bibitem{ceretani} A. N. Ceretani, N. N. Salva and D. A. Tarzia.
An exact solution to a Stefan problem with variable thermal conductivity and
a Robin boundary condition (2017). \url{https://arxiv.org/abs/1706.06984}.
\bibitem{code} Transport Verification Lab. Versioned source, synthetic
experiments and validation records.
\url{https://github.com/racostacme-come/transport-verification-lab}.
The manuscript build report pins source and result SHA-256 hashes.
\end{thebibliography}
\end{document}
